\section{Conclusion}\label{sec:conclusion}

We set out to test whether US aid, instrumented by UN Security Council membership, induces countries to express greater pro-US sentiment in their UN General Debate speeches. Our findings are consistently null. The UNSC instrument fails the first stage under every fixed-effect specification, including the region--year FE design of Kuziemko and Werker (2006). The reduced form is negative across specifications---UNSC members are, if anything, slightly \emph{less} pro-US in their rhetoric---though not statistically significant at conventional levels. An event study reveals no systematic dynamic pattern around the UNSC term beyond a marginally significant lead-2 coefficient, potentially reflecting pre-election diplomatic activity.

These results contribute three lessons. First, they demonstrate the fragility of the UNSC--aid IV chain: what appears as a significant relationship in simpler specifications disappears when country fixed effects absorb between-country variation. Second, they highlight the gap between diplomatic rhetoric and observable behavior: our LLM-based stance measure, despite achieving acceptable inter-coder agreement, barely correlates with UN voting patterns. Third, they illustrate the value of pre-registered causal gates: by committing to the $F>10$ threshold before inspecting results, we avoid the temptation to report weak-instrument 2SLS estimates.

The null result is not a disappointment. It is evidence that, under more stringent identification than the existing literature, the material-power-buys-allegiance thesis finds no empirical support through the UNSC--aid channel. Future work should explore alternative instruments for aid, alternative measures of expressed allegiance, and alternative forums where diplomatic speech may be more costly and therefore more informative about genuine preferences.
